\subsection{Training-data and selection factors}
A retrospective $2\times2$ ablation crosses original/controlled training scenes with F1/balanced-accuracy clip selection for mean and attention pooling. Twelve new and twelve retained runs cover eight conditions with three seeds each, keeping frame evaluation fixed. Under balanced selection, attention's controlled-test frame BA is \OriginalAttentionBA{} with original training data and \ControlledAttentionBABA{} with controlled training. Controlled training under F1 selection reaches \ControlledAttentionFoneBA{}. The supplement gives paired-seed contrasts. This separates two protocol factors, but the data factor still includes extra scenes and training exposure, and these test tracks were already inspected.

\subsection{Track stability and conditional errors}
Retrained attention exceeds mean on each of the four test tracks. The unweighted controlled-track BA difference averages \TrackDelta{}; leaving out any one track gives \TrackLooLow{}--\TrackLooHigh{}. This checks sensitivity to the sampled tracks without treating clips or seeds as independent musical works; it is not a population confidence interval.

With bass absent, bass FPR averaged over the three attention seeds is \BassFprDrumsOn{} when drums are active (1235 reference bins) and \BassFprDrumsOff{} when drums are absent (1055 bins). Thus co-occurrence is associated with a measurable error difference, beyond the overlap of clip labels alone. The supplement supplies conditional contrasts and denominators; causal interpretation is considered in the discussion.
